\documentclass{beamer}
\usepackage{multirow}
\usepackage{listings}

\usepackage{tikz}
\usetikzlibrary{positioning}
\usetikzlibrary{calc}
\usetikzlibrary{shadows}
\usetikzlibrary{fit}

\mode<presentation>
{
\usetheme{boxes}
\usecolortheme{crane}
%\usetheme{Warsaw}
\setbeamercovered{transparent}
}


\usepackage[polish]{babel}
\usepackage[T1]{fontenc}
\usepackage[utf8]{inputenc}
\usepackage{psfrag}
\usepackage{graphicx}
\usepackage{times}

\usepackage{listings}


\makeatletter
\newenvironment{minsizebox}[2]{%
    \pgfmathsetlength\@tempdima{#2}%
    \pgfmathsetlength\pgf@yc{\pgfkeysvalueof{/pgf/inner ysep}}%
    \advance\@tempdima by -2\pgf@yc
    \begin{lrbox}{\@tempboxa}%
        \begin{minipage}[t]{#1}%
           \vspace{0pt}%
}{%
        \end{minipage}%
    \end{lrbox}%
    \ifdim\@tempdima>\dp\@tempboxa
        \dp\@tempboxa=\@tempdima
    \fi
    \box\@tempboxa
}
\makeatother
\tikzset{onslide/.code args={<#1>#2}{%
  \only<#1>{\pgfkeysalso{#2}} % \pgfkeysalso doesn't change the path
}}

\newenvironment{pic}[1]{%
\begin{center}\begin{tikzpicture}#1}{%
\end{tikzpicture}
\end{center}
\end{lstlisting}}


\newcommand{\listing}{\node[listing]{\begin{lstlisting}
\end{lstlisting}
}
}


\tikzset{common/.style={fill=white,draw=black,rounded corners=1mm}}
\tikzset{listing/.style={fill=white,draw=black,font=\scriptsize,
rounded corners=0.5mm,drop shadow}}
\tikzset{title/.style={fill=white,font=\Large,rectangle, draw=black,
drop shadow, minimum width=\textheight, minimum height=15mm,
 rounded corners =3mm}}

\tikzset{display/.style={fill=white,rectangle, draw=black,
drop shadow, minimum width=8cm, minimum height=15mm,
 rounded corners =3mm}}

\renewcommand{\v}[1]{\vec{#1}}
\newcommand{\av}[1]{\langle#1\rangle}
\renewcommand{\H}{\mathcal{H}}
\renewcommand{\a}[1]{\v{a}_{#1}}
\renewcommand{\L}{\mathcal{L}}

\newcommand{\vphi}{\vec{\varphi}}
\newcommand{\vi}{{\vec{i}}}
\newcommand{\vj}{{\vec{j}}}
\newcommand{\vmu}{\vec{\mu}}

\newcommand{\redblack}[3]{ 
\pgfmathtruncatemacro{\odd}{Mod(#1+#2,2)}
   \ifnum \odd=0
      \fill[fill=black] (#1,#2) circle[radius=#3];
   \else
      \fill[fill=red] (#1,#2) circle[radius=#3];
   \fi
}

\newcommand{\fillblockgray}[2]{
\pgfmathtruncatemacro\llx{\bksize*(#1)-2}
\pgfmathtruncatemacro\lly{\bksize*(#2)-2}
\pgfmathtruncatemacro\urx{\bksize*(#1+1)+1}
\pgfmathtruncatemacro\ury{\bksize*(#2+1)+1}
\foreach \x in {\llx, ..., \urx}
  \foreach \y in {\lly,...,\ury} {
    \fill[craneorange] (\x, \y) circle(.25);
  }
}

\newcommand{\npartition}[3]{
\pgfmathtruncatemacro\row{Mod(#2,4)}
\ifnum \row = 0
   \pgfmathtruncatemacro\num{Mod(#1,4)}
\else \ifnum \row  = 1
   \pgfmathtruncatemacro\num{Mod(#1,4)+4}
\else \ifnum \row  = 2
   \pgfmathtruncatemacro\num{Mod(Mod(#1,4)+2,4)}
\else \ifnum \row  = 3
   \pgfmathtruncatemacro\num{Mod(Mod(#1,4)+2,4)+4}
\fi
\fi
\fi
\fi
\ifnum \num = #3
\fill[black] (#1, #2) circle(0.35);
\fi

}

\newcommand{\markpartition}[3]{
\pgfmathtruncatemacro\llx{\bksize*(#1)}
\pgfmathtruncatemacro\lly{\bksize*(#2)}
\pgfmathtruncatemacro\urx{\bksize*(#1+1)-1}
\pgfmathtruncatemacro\ury{\bksize*(#2+1)-1}
\foreach \x in {\llx, ..., \urx}
  \foreach \y in {\lly,...,\ury} {
    \npartition{\x}{\y}{#3}
  }
}


\newcommand{\partition}[3]{
\pgfmathtruncatemacro\lchalf{\bksize / 2}
      \pgfmathtruncatemacro\my{Mod(#1, 2)}
      \ifnum \my = 0
        \pgfmathtruncatemacro\ny{Mod(#1, 4)}
        \ifnum \ny = 2
          \pgfmathtruncatemacro\mx{Mod(#2+2, 4)}
        \else
          \pgfmathtruncatemacro\mx{Mod(#2, 4)}
        \fi
        \ifnum \mx = 0
            \fill[red] (#2, #1) circle(#3);
          \fi
         \fi         
       }


\definecolor{colorone}{HTML}{006633} % default 116699
\definecolor{colortwo}{HTML}{DDFF99} % default CCCCCC
\def\bksize{8}
\def\bkcount{4}
\def\lcsize{5}

\lstset{basicstyle=\small}


\title{GPU implementation of the Monte-Carlo simulations of the extended Ginzburg-Landau model.}
\date{10 September 2013}
\author{Piotr Bialas}
\institute{Faculty of Physics, Astronomy and Computer Science, Jagiellonian University \\
Department of Game Technology \\
\vspace{1em}
\includegraphics[height=0.18\paperheight]{../../UJ-logo}}
\begin{document}

\begin{frame}
\titlepage
\vfill
J.~Kowal, A.~Strzelecki
\vfill
\end{frame}



\begin{frame}
\frametitle{Outline}
\tableofcontents
\end{frame}

\section{The model}


\begin{frame}
\frametitle{Ginsburg-Landau model ($\vphi^4$)}
\begin{equation*}\label{eq:ham}\begin{split}
\H[\varphi]&=\sum_{\vi}\Biggl(
\frac{1}{2}\sum_{\mu=1}^d(\vphi(x_\vi+\a{\mu})-\vphi(x_\vi))^2\\
&\phantom{=\sum_{\vi}\bigl(}+\frac{\mu^2}{2}|\vphi(x_{\vi})|^2+
\frac{g}{24}(|\vphi(x_{\vi})|^2)^2\\
&\phantom{=\sum_{\vi}\bigl(}+
\frac{1}{2\Lambda}\Bigl(\sum_{\mu=1}^d
(\vphi(x_\vi+\a{\mu})-2\vphi(x_{\vi})+\vphi(x_\vi-\a\mu))\Bigr)^2\Biggr).
\end{split}
\end{equation*}
\begin{equation*}
P[\varphi]\propto \exp\left(-H[\varphi]\right)
\end{equation*}
\end{frame}

\begin{frame}
\frametitle{Monte-Carlo}
\[
\av{O}=\mathcal{Z}^{-1}\int\mathcal{D}\phi 
O\bigl(\v\phi(x)\bigr)
\exp
\biggl(
-\mathcal{H}\bigl[\v\phi(x)\bigr]
\biggr)
\]


\[
\av{O}\approx\frac{1}{N}\sum O(\v\phi_i)
\]

$\phi_i$ idenpendent samples drawn from  distribution

\[
\mathcal{Z}^{-1}
\exp
\biggl(
-\mathcal{H}\bigl[\v\phi\bigr]
\biggr)
\]

\end{frame}


\begin{frame}
\frametitle{Metropolis}

\begin{align*}
\phi_{new}(x_i)&=\phi(x_i)+\epsilon\\
\Delta\L&=\H[\v\phi_{new}]-\H[\v\phi]
\end{align*}
We assign  $\phi(x_i)=\phi_{new}(x_i)$ with probability
\[
\min\bigl(1,\exp(-\Delta\H)\bigr)
\]

\[
r < exp(-\Delta\H) ,\qquad r \in [0,1)  \text{ random}
\]
\end{frame}


\begin{frame}
\frametitle{Ginsburg-Landau model ($\vphi^4$)}
\begin{tikzpicture} 
%\pgfmathtruncatemacro\lcsize{\bksize * \bkcount}
\def\lcsize{8}
\pgfmathtruncatemacro\lcsizem{\lcsize - 1}
\pgfmathtruncatemacro\lcsizeb{\lcsize - \bksize}
\pgfmathtruncatemacro\bksizea{\bksize - 2}
\pgfmathtruncatemacro\bksizeb{\bksize - 1}
\pgfmathtruncatemacro\bksizec{\bksize * 2}
\pgfmathtruncatemacro\bksizem{\bksizec - 1}
\pgfmathtruncatemacro\bksized{\bksizec + 1}
\draw[very thin, gray] (-.5, -.5) grid (\bksize - .5, \bksize - .5);
\foreach \x in {0,...,\lcsizem} 
\foreach \y in {0, ...,\lcsizem}
\draw[fill=white,draw=black]  (\x, \y) circle[radius=0.1];

\coordinate (c) at (4,4);

\node at (c)  [circle,fill=red,draw=black] {};

\def\coronaa{(0,1), (0,-1), (1,0), (-1,0)}
\def\coronab{(0,2), (0,-2), (2,0), (-2,0)}
\def\coronac{(1,1), (1,-1), (-1,-1), (-1,1)}

\foreach \p in \coronaa {
\path (c) +\p node  [circle,fill=gray]    {};
}

\foreach \p in \coronab {
\path (c) +\p node  [circle,fill=gray]    {};
}

\foreach \p in \coronac {
\path (c) +\p node  [circle,fill=gray]    {};
}


\end{tikzpicture}

\end{frame}

\begin{frame}
\frametitle{Parallel partitions}
\begin{tikzpicture}[scale=0.75]

% \clip[rounded corners=0cm] (-1, -1) rectangle (\bksize + 1, \bksize + 1);

% \draw[help lines] (-1, -1) grid (\bksize + 1, \bksize + 1);

\pgfmathtruncatemacro\bksizem{\bksize - 1}

% % lattice points
 \foreach \x in {0, ..., \bksizem}
   \foreach \y in {0, ..., \bksizem} {
     \pgfmathtruncatemacro\four{Mod(\y,4)}
     \pgfmathtruncatemacro\eight{Mod(\x,4)+4*Mod(\y,2)}
     \ifnum \four<2
     \pgfmathtruncatemacro\eight{Mod(\x,4)+4*Mod(\y,2)}
     \node[draw,rectangle,minimum size=0.6cm] at (\x,\y)  {\eight};
     \else
     \pgfmathtruncatemacro\eight{Mod(\x+2,4)+4*Mod(\y,2)}
     \node[draw,rectangle,minimum size=0.6cm] at (\x,\y)  {\eight};
     \fi
}
% neighborhood
\def\neighborhood{--
  ++( 0,  -.5)-- ++( 1,  0)-- ++( 0, -.5 )-- ++(0, -.5)-- ++( 1, 0  )--
  ++( 0, -1.0)-- ++( 1,  0)-- ++( 0,  1.0)-- ++(1, 0  )-- ++( 0, 1.0)--
  ++( 1,  0  )-- ++( 0,  1)-- ++(-1,  0  )-- ++(0, 1  )-- ++(-1, 0  )--
  ++( 0,  1  )-- ++(-1,  0)-- ++( 0, -1  )--
  ++(-1,  0  )-- ++( 0, -1)-- ++(-1,  0  )-- ++(0, -.6)}

\pgfmathtruncatemacro\lchalf{\bksize / 2}


% main center neighborhood
\path[draw=red,  rounded corners=.1cm, shorten >=2pt, line width=.75mm]
(\lchalf - 2.5, \lchalf) \neighborhood;

\end{tikzpicture}

\end{frame}

\begin{frame}
\frametitle{Random number generation}
\begin{itemize}
\item Generating random number involves changing state.
\item Two processes cannot do it at the same time.  
\item Every process/thread  must have his own random number generator. 
\item We use Hybrid Tausworthe generator from GPU GEMS III.
\end{itemize}
\end{frame}


\section{CUDA implementation}




\begin{frame}
\frametitle{Memory hierarchy}
\begin{itemize}
\item Max throughput
\begin{itemize}
\item 32 FMAD per cycle
\begin{itemize}
\item 12 bytes in 4 bytes out
\item 12x 32x15x1401MHz
\item 8.07 TB/s in 2.67 TB/s out
\end{itemize}
\end{itemize}
\item Registers can sustain it.
\item Shared memory
\begin{itemize}
\item 4 x 32 x 15 x 1401MHz/2 = 1.3TB/s
\end{itemize}
\item Global memory
\begin{itemize}
\item 156GB/s
\end{itemize}
\end{itemize}

\end{frame}

\begin{frame}

\begin{tikzpicture}[scale=0.25]

\pgfmathtruncatemacro\lcsize{\bksize * \bkcount}
\pgfmathtruncatemacro\lcsizem{\lcsize - 1}
\pgfmathtruncatemacro\lcsizeb{\lcsize - \bksize}


\draw[very thin, gray] (-.5, -.5) grid (\lcsize - .5, \lcsize - .5);

\foreach \x in {0, \bksize, ..., \lcsizeb}
  \foreach \y in {0, \bksize, ..., \lcsizeb} {
    \pgfmathtruncatemacro\n{Mod(\x/\bksize, 2)+2 * Mod(\y/\bksize, 2)+1}
    \draw[colorone, xshift=\x cm, yshift=\y cm, rounded corners=2pt]
      (-.4, -.4) rectangle
      node[red, fill opacity=0, anchor=center, font=\Huge\bfseries\sffamily]
      {\n} (\bksize - .6, \bksize - .6);
  }

\foreach \x in {0, ..., \lcsizem}
  \foreach \y in {0, ..., \lcsizem}{
 
    \path[draw=black, fill=white] (\x, \y) circle(.25);
}


\only<1>{\fillblockgray{0}{0}
\markpartition{0}{0}{0}
\fillblockgray{0}{2}
\markpartition{0}{2}{0}
\fillblockgray{2}{0}
\markpartition{2}{0}{0}
\fillblockgray{2}{2}
\markpartition{2}{2}{0}}

\only<2>{\fillblockgray{0}{0}
\markpartition{0}{0}{1}
\fillblockgray{0}{2}
\markpartition{0}{2}{1}
\fillblockgray{2}{0}
\markpartition{2}{0}{1}
\fillblockgray{2}{2}
\markpartition{2}{2}{1}}

\only<3>{\fillblockgray{0}{0}
\markpartition{0}{0}{2}
\fillblockgray{0}{2}
\markpartition{0}{2}{2}
\fillblockgray{2}{0}
\markpartition{2}{0}{2}
\fillblockgray{2}{2}
\markpartition{2}{2}{2}}

\only<4>{
\fillblockgray{1}{0}
\markpartition{1}{0}{0}
\fillblockgray{1}{2}
\markpartition{1}{2}{0}
\fillblockgray{3}{0}
\markpartition{3}{0}{0}
\fillblockgray{3}{2}
\markpartition{3}{2}{0}}

\only<5>{
\fillblockgray{1}{0}
\markpartition{1}{0}{1}
\fillblockgray{1}{2}
\markpartition{1}{2}{1}
\fillblockgray{3}{0}
\markpartition{3}{0}{1}
\fillblockgray{3}{2}
\markpartition{3}{2}{1}}


\only<6>{
\fillblockgray{1}{0}
\markpartition{1}{0}{2}
\fillblockgray{1}{2}
\markpartition{1}{2}{2}
\fillblockgray{3}{0}
\markpartition{3}{0}{2}
\fillblockgray{3}{2}
\markpartition{3}{2}{2}}


\end{tikzpicture}

\visible<1>{\tiny M.~Weigel, J.~Comput.~Phys. 231 (2012) 3064.}
\end{frame}


\section{CPU implementation}

\begin{frame}
\frametitle{CPU}
\begin{itemize}
\item First measurements shoved over 300 fold  performance increase over CPU!
\item Not fair.
\item Compared to a non--parallel, non-vectorized program.
\end{itemize}
\end{frame}



\begin{frame}
\frametitle{CPU implementation}
\begin{itemize}
\item We use same structure as on GPU.
\item We use OpenMP to distribute lattice blocks across cores.
\item In each block we use SSE instructions to process 4 or 8 lattice points in parallel. 
\end{itemize}
\end{frame}

\begin{frame}

\begin{tikzpicture}[scale=0.25]

\pgfmathtruncatemacro\lcsize{\bksize * \bkcount}
\pgfmathtruncatemacro\lcsizem{\lcsize - 1}
\pgfmathtruncatemacro\lcsizeb{\lcsize - \bksize}


\draw[very thin, gray] (-.5, -.5) grid (\lcsize - .5, \lcsize - .5);

\foreach \x in {0, \bksize, ..., \lcsizeb}
  \foreach \y in {0, \bksize, ..., \lcsizeb} {
    \pgfmathtruncatemacro\n{Mod(\x/\bksize, 2)+2 * Mod(\y/\bksize, 2)+1}
    \draw[colorone, xshift=\x cm, yshift=\y cm, rounded corners=2pt]
      (-.4, -.4) rectangle
      node[red, fill opacity=0, anchor=center, font=\Huge\bfseries\sffamily]
      {\n} (\bksize - .6, \bksize - .6);
  }

\foreach \x in {0, ..., \lcsizem}
  \foreach \y in {0, ..., \lcsizem}{
 
    \path[draw=black, fill=white] (\x, \y) circle(.25);
}


\only<1>{\fillblockgray{0}{0}
\markpartition{0}{0}{0}
\fillblockgray{0}{2}
\markpartition{0}{2}{0}
\fillblockgray{2}{0}
\markpartition{2}{0}{0}
\fillblockgray{2}{2}
\markpartition{2}{2}{0}}
\end{tikzpicture}
\end{frame}


%%%%%%%%%%%%%%%%%%%%%%%%%%%%%%%%%%%%%%%%%%%%%%%%%%%%%%%%%%%%%%%%%%%%%%%%%%%%%%%%

\begin{frame}
\frametitle{SSE/AVX}
\begin{itemize}
\item Modern Intel and AMD processors have vector operations.
\begin{itemize} 
\item Streaming Side Extension  (SSE) Four single precision operations.
\item Advanced Vector Extensions (AVX) Eight  single precision operations.
\item Xeon $\phi$ 16  singleprecision operations.
\end{itemize}
\item Not using vector operations wastes up to 7/8 of CPU power!.  
\end{itemize}
\end{frame}

\begin{frame}
\frametitle{SSE/AVX}
\begin{itemize}
\item Programming  with vector operations can be done
\begin{itemize}
\item Asembler :)
\item Instrinsincs
\item Vector types
\item Compiler directives (Intel 13)
\end{itemize}
\end{itemize}
\end{frame}



\begin{frame}
  \frametitle{AVX/SSE  library}
  \begin{itemize}
    \item Custom class wrapping \texttt{vector\_size}
    \item Drop in vector replacement for scalar type
    \item Defines all operators available for scalar types
    \item Existing function just take vector arguments now
    \item Rest remains the same \\
    \texttt{action -= phi2*(quadratic\_coef+gr*phi2)}
  \end{itemize}
\end{frame}

%%%%%%%%%%%%%%%%%%%%%%%%%%%%%%%%%%%%%%%%%%%%%%%%%%%%%%%%%%%%%%%%%%%%%%%%%%%%%%%%

\begin{frame} 

\frametitle{Performance}

\begin{center}
\def\arraystretch{1.1}%
\setlength{\tabcolsep}{0.25em}
\begin{tabular}{ r| r r r | r r r| r }
           & \multicolumn{3}{c|}{GPU\footnotemark} &
             \multicolumn{3}{c|}{CPU\footnotemark}  &  \\
    size   & time(ns) & Gflops & Gexps &time(ns) & Gflops & Glogs & boost \\
  \hline
   $32^3$  & 1.26     &  27.1 & 0.8 & 3.88    &  8.8 & 0.26 &  $3.1 \times$ \\
   $64^3$  & 0.17     & 206.8 & 6.1 & 3.09    & 11.0 & 0.32 & $18.8 \times$ \\
  $128^3$  & 0.13     & 273.0 & 8.0 & 3.31    & 10.3 & 0.30 & $26.5 \times$ \\
  $256^3$  & 0.11     & 302.0 & 8.9 & 3.33    & 10.3 & 0.30 & $29.5 \times$ \\
  \hline
                      & \footnotesize peak & $\approx 1000$ &
                      & \footnotesize peak & $\approx 160$ &
\end{tabular}
\end{center}
\vfill
\small
$^1$\emph{GeForce GTX 470} \, 1.2 GHz, 1.25GB, 448 CUDA cores \\
$^2$\emph{Intel Core i5-2400S} \, 2.5 GHz, 4 core, 6MB cache, SSE4.x/AVX
\end{frame}

%%%%%%%%%%%%%%%%%%%%%%%%%%%%%%%%%%%%%%%%%%%%%%%%%%%%%%%%%%%%%%%%%%%%%%%%%%%%%%%%

\section{Conclusions}

\begin{frame}
  \frametitle{Conclusions}
  \begin{itemize}
    \item GPU has more computing power.
    \item GPU is more efficient: easier to reach closer to the peak.
    \begin{itemize}
      \item Possible thanks to the use of fast shared memory
      \item Using cache is just not enough
      \item Shared memory simplifies indexing
    \end{itemize}
    \item The CUDA (SIMT) execution model can be ported to CPU:
    \begin{itemize}
      \item No shared memory, only cache, little control
      \item No gather/scatter vector load operations
      \item Vector assembly from memory is a bottleneck
      \item Half-wide only (128-bit) vector integer operations
      \vspace{1em}
      \item \emph{Xeon $\phi$} \&  \emph{Haswell} address (some) these issues
    \end{itemize}
  \end{itemize}
\end{frame}

\end{document}
